\begin{enumerate}
	\modo{}{\color{white}}

	% subtração não-negativa 
	\item Crie um programa que exiba a subtração de dois números, substituindo-a por 0 caso ela for negativa.

	% divisor diferente de 0
	\item Crie um programa que exiba a divisão de dois números caso o divisor for diferente de 0 e exiba uma mensagem de erro caso não exista resultado.
	
	Altere o programa para a mensagem de erro ser diferente para $\frac{0}{0}$, informando que o resultado é indeterminado. 

\quebra

	% graus, grados e radianos 
	% agudo, reto, obtuso, raso 
	\item Crie um programa que faça a conversão de um ângulo escrito pelo usuário. 
	Receba do usuário qual unidade de medida foi inserida na entrada (graus ou radianos)
	e exiba-a convertida para a outra unidade.
	
	Tente alterar o programa para incluir também entrada ou saída em grados.
	
	\begin{multicols}{2}	
		\begin{equation*}
			\frac{10}{9} \text{gr} = 1^\circ
		\end{equation*}
		
		\begin{equation*}
			1 \text{gr} = 0,9^\circ
		\end{equation*}
	\end{multicols}
	
	Inclua também no programa a exibição de se o ângulo inserido é agudo ($< 90^\circ$), reto ($= 90^\circ$), obtuso ($> 90^\circ$ e $< 180^\circ$), raso ($= 180^\circ$), reflexo ($> 180^\circ$ e $< 360^\circ$) ou completo ($= 360^\circ$).
	
	% Celsius, Fahrenheit e Kelvin  
	% água está sólida, líquida ou gasosa?
	\item Crie um programa que faça a conversão de uma temperatura escrita pelo usuário.
	Receba do usuário qual unidade de medida foi inserida na entrada (graus Celsius, Fahrenheit ou Kelvin)
	e a unidade para converter e a exiba.	
	
	\begin{multicols}{2}	
		\begin{equation*}
			C = K - 273,16
		\end{equation*}
		
		\begin{equation*}
			K = C + 273,16
		\end{equation*}
		
		\begin{equation*}
			F = 1,8(K - 273,16) + 32
		\end{equation*}
		
		\begin{equation*}
			K = \frac{F-32}{1,8} + 273,16 
		\end{equation*}
	\end{multicols}
	
	Inclua também no programa a exibição de se a água destilada a 1 atm (pressão do nível do mar) estaria sólida ($< 0^\circ C$), líquida ($> 0^\circ C$ e $< 100^\circ C$), gasosa ($> 100^\circ C$) ou se estaria em algum dos pontos de transição de estado.  	

	% Ano bissexto 
	% juliano 
	% gregoriano 
	\item Crie um programa que receba um ano do usuário e informe se ele é bissexto ou não.
	
	Anos múltiplos de 4 são bissextos no calendário juliano.
	
	Anos múltiplos de 400 ou de 4, mas não de 100, são bissextos no calendário gregoriano.
	
	Exiba se o ano é bissexto para cada calendário.
	% repetido 
	
	
\textbf{Altere os exercícios anteriores para repetir as ações várias vezes, aceitando entradas e exibindo resultados até alguma condição de parada ser satisfeita (crie algum critério e informe o usuário)}

\quebra	
	
	\item Crie um programa que exiba cada linha de texto escrita pelo usuário até que seja inserida uma linha vazia.
	
	\item Crie um programa que receba vários números, até que seja inserido 0, e exiba a soma deles.
	
	\item Crie um programa que exiba a soma de dois números escritos pelo usuário, repetindo a entrada de dados e o cálculo da soma até que os dois números inseridos sejam 0.
	
	Altere o programa para parar quando o usuário inserir uma linha vazia, sem nenhum número.		
	
\quebra	

	% distância entre pontos com d dimensões 
	\item Faça um programa que exiba a distância euclidiana entre dois pontos definidos no código.
	
	\item Faça um programa que receba das coordenadas de dois pontos e exiba a distância euclidiana entre eles. 
	
	Defina as dimensões (quantidades de eixos) no código, permitindo alteração fácil.	

\quebra	
	
	% heron com vetores 
	\item Faça um programa que receba os comprimentos dos lados $l_1$, $l_2$ e $l_3$ de um triângulo e exiba a área de acordo com o Teorema de Heron
	
	\begin{equation*}
		A = \sqrt{p(p-l_1)(p-l_2)(p-l_3)}
	\end{equation*}	
	
	Onde $p$ é o semiperímetro:
	
	\begin{equation*}
		p = \frac{l_1+l_2+l_3}{2}
	\end{equation*}
	
	Use laços de repetição e listas para diminuir o código.
	
\quebra	
	
	% é um triângulo mesmo?
	\item Faça um programa que receba os comprimentos dos lados $a$, $b$ e $c$ e informe se eles podem construir um triângulo plano.
	
	A condição de existência do triângulo é $a < b + c$, $b < a + c$ e $c < a + b$.
	
	\item Faça um programa que receba os comprimentos dos lados $a$, $b$ e $c$ de um triângulo e informe se ele é 
	equilátero (todos os comprimentos são iguais), 
	isósceles (somente dois lados têm o mesmo comprimento)
	ou escaleno (todos os lados são diferentes).
	
	\item Faça um outro programa que receba as medidas dos vários lados de uma figura geométrica e informe se ela é um polígono regular (todos os lados têm o mesmo comprimento) ou não.
	
	Faça de duas maneiras: 
	a primeira recebendo todos os lados até ser inserido 0 
	e a segunda recebendo a quantidade de lados primeiro e depois recebendo todos os lados, parando ao ter a quantidade informada. 
	 


	% bhaskara se for quadrática, linear se não for	
	\item Faça um programa que receba os coeficientes $a$, $b$ e $c$ de uma equação $ax^2 + bx + c = 0$ e exiba as raízes $x_1$ e $x_2$ se a equação for quadrática ou a raiz $x$ se for uma equação linear.
	
	Caso a equação quadrática não tiver solução real ($\Delta < 0$), exiba uma mensagem de erro. 

	Também exiba informações de erro se a equação linear não tiver solução única ($b = c = 0$) ou se não tiver solução ($b = 0$ e $c \ne 0$).

	% Sistema linear genérico 
	\item Faça um programa que receba os coeficientes $A_{11}$, $A_{21}$, $A_{22}$, $A_{31}$, $A_{32}$, $A_{33}$ .... $A_{i1}$ ... $A_{ii}$, $b_1$, $b_2$, $b_3$ ... $b_i$ de um sistema de equações lineares escalonado e exiba os valores de $x_1$, $x_2$, $x_3$ ... $x_i$.		

	Defina $i$ no código ou receba do usuário antes da leitura.	 

\quebra

	% Números de 1 a 10
	\item Crie um programa que exiba os inteiros de 1 a 10 (incluindo ambos). 

	% Números de 1 a n
	\item Crie um programa que exiba os inteiros de 1 a $n$.

	Faça versões com e sem entrada de dados

	% Tabuada 
	\item Crie um programa que exiba a tabuada de $k$.

	Altere o programa para exibir $n$ linhas da tabuada.  

	% Triângulo de Pascal 
	\item Crie um programa que exiba as primeiras $n$ linhas do Triângulo de Pascal. 
	
	\textbf{Dica}:
	para exibir um valor sem quebrar a linha ao final, use a opção \texttt{sep}
	
	\begin{verbatim}
		print('Não será quebrada', sep='')
	\end{verbatim}

	% Fatorial 
	\item Crie um programa que calcule o fatorial de $n$.
	
	\begin{equation*}
		n! = 1 \times 2 \times 3 \times .... \times n
	\end{equation*}	
	
	Lembre-se que $0! = 1$.
	Exiba uma mensagem de erro caso $n < 0$.
	

	\item Crie um programa que calcule o arranjo simples.
	
	\begin{equation*}
		A^k_n = \frac{n!}{(n-k)!} = (n-k+1) (n-k+2) (n-k+3) .... n
	\end{equation*}

	% Número binomial 
	\item Crie um programa que calcule o número binomial de $n$ na classe $k$.
	
	\begin{equation*}
		\binom{n}{k} = C^k_n = \frac{n!}{k!(n-k)!}
	\end{equation*} 
	
	Exiba uma mensagem de erro caso $n < 0$, $k < 0$ ou $n < k$.

\quebra 	

	% Fibonacci 
	\item Faça um programa que exiba os primeiros $n$ termos da sequência de Fibonacci. 
	Os primeiros termos são $F_0 = 0$ e $F_1 = 1$ e cada termo seguinte é a soma dos dois anteriores.
	
	\begin{equation*}
		F_n = F_{n-1} + F_{n-2}
	\end{equation*}
	
	Altere o programa para aceitar $n < 0$ também.
	
	\begin{equation*}
		F_{n+2} - F_{n+1} = F_n 
	\end{equation*}		
	
	Tente também fazer essa alteração com somente um laço de repetição e executando o condicional uma vez.
	
	\item Faça um programa que exiba os primeiros $n$ termos de uma sequência parecida com a Fibonacci.
	Receba os $k$ termos iniciais e calcule cada termo seguinte como a soma dos $k$ termos anteriores. 
	
	\begin{equation*}
		T_n = T_{n-1} + T_{n-2} + T_{n-3} + ... + T_{n-k+1} + T_{n-k}
	\end{equation*}
	
	Altere o programa para aceitar $n < 0$ também.
	
	\begin{equation*}
		T_{n+k} - T_{n+k-1} - T_{n+k-2} - T_{n+k-3} - ... - T_{n+1} = T_n 
	\end{equation*}		
	
	 
	\item Crie um programa que receba um número inteiro $p$ e exiba se ele é primo, composto ou nenhum.
	
	Um inteiro $p$ é primo se ele só for divisível por 1, -1, $p$ e $-p$ (sendo esses quatro distintos). 	
	Um inteiro é composto se for divisível por mais do que estes quatro números.
	0, -1 e 1 não são nem compostos, nem primos.
	

	

\quebra

	% 27 PA
	\item Crie um programa que 
	receba os valores de $a_1$, $r$ e $n$ 
	e exiba os primeiros $n$ termos de uma 
	Progressão Aritmética (PA).

	\begin{equation*}			
		a_n = a_{n-1} + r
	\end{equation*}	

	Faça a soma termo a termo e também calcule com a fórmula:

	\begin{equation*}
		S_n = n\frac{a_1 + a_n}{2} = n\frac{2 a_1 + r(n - 1)}{2}
	\end{equation*}

	Exiba a soma feita das duas formas.		

	% 28 PG
	\item Crie um programa que 
	receba os valores de $a_1$, $q$ e $n$
	e exiba os primeiros $n$ termos de uma 
	Progressão Geométrica (PG).

	\begin{equation*}
		a_n = a_{n-1} q
	\end{equation*}	

	Faça a soma termo a termo e também calcule com a fórmula:

	\begin{equation*}
		S_n = \frac{a_n q - a_1}{q - 1} = a_1 \frac{q^n - 1}{q - 1}
	\end{equation*}

	Exiba a soma feita das duas formas.


	% Juros simples 
	\item Crie um programa que receba o valor do capital $C$, da taxa de juros simples $i$, o tempo $t$ e exiba o montante $M_t$ de acordo com a fórmula:
	\begin{multicols}{2}
		\begin{equation*}
			J_t = C i t
		\end{equation*}	

		\begin{equation*}
			M_t = C + J_t
		\end{equation*}	
	\end{multicols}
	Altere o programa para calcular com a fórmula:		
	\begin{equation*}
		M_t = C (1 + i t) 
	\end{equation*}
	
	Altere também o programa para calcular o montante período a período 
	\begin{multicols}{2}
		\begin{equation*}
			M_0 = C
		\end{equation*}
		
		\begin{equation*}
			M_t = M_{t-1} + C i
		\end{equation*}
	\end{multicols}	
	% com parcelas		
	Altere o programa para, a cada termo, receber do usuário o pagamento de uma parcela e descontar do montante, encerrando o programa quando o valor for quitado.		

\quebra	 

	% Juros compostos 
	\item Crie um programa que receba o valor do capital $C$, da taxa de juros compostos $i$, o tempo $t$ e exiba o montante $M_t$ de acordo com a fórmula:
	\begin{equation*}
		M_t = C (1 + i)^t
	\end{equation*}	
	Altere também o programa para calcular o montante período a período 
	\begin{multicols}{2}	
		\begin{equation*}
			M_0 = C
		\end{equation*}		

		\begin{equation*}
			M_t = M_{t-1} (1 + i)
		\end{equation*}
	\end{multicols}
	% com parcelas		
	Altere o programa para, a cada termo, receber do usuário o pagamento de uma parcela e descontar do montante, encerrando o programa quando o valor for quitado. 

	% Média
	\item Crie um programa que receba as notas das avaliações 1, 2, 3.... $n$ e seus pesos e exiba a média ponderada delas.

	% Votos válidos 
	\item Faça um programa que receba a quantidade de votos dos candidatos 1, 2, 3.... $n$, de votos brancos e votos nulos e exiba as porcentagens de votos válidos de cada um dos candidatos.			 

	% Receita 
	\item Faça um programa que receba as quantidades de $n$ ingredientes e seus nomes, o tamanho da porção da receita e o tamanho da porção desejada e exiba os nomes e as quantidades necessárias para a porção esperada.	

	% Matrículas e notas 
	\item Faça um programa em que estejam registrados os números de matrícula e os nomes dos estudantes de uma turma. 
	
	Receba uma matrículas do usuário.
	Para cada matrícula, exiba o nome do estudante, se a matrícula já estiver registrada. 
	
	Caso a matrícula não estiver registrada, receba o nome do estudante e registre-o.
	
	Altere o programa para registrar as notas também.
	Receba do usuário se ele quer inserir mais notas numa matrícula, editar o nome ou exibir as informações já registradas.
	
	\item Faça um programa que receba os votos de candidatos registrados com número e nome. 
	Defina os nomes e números dos candidatos no código.
	
	Cada número recebido deve ser contabilizado para o candidato registrado e exibir seu nome.
	
	Se o número não tiver candidato registrado, considere e exiba voto nulo.
	
	Se o voto for 0, considere e exiba voto branco.
	
	Pare de receber votos quando a quantidade máxima $v$ de votantes da seção (número definido no código) 
	for atingido ou quando for inserido número negativo (caso a seção seja fechada com abstenções ao fim do horário).
	
	Ao encerrar a votação, exiba as quantidades de votos de cada candidato, brancos e nulos
	 e calcule as porcentagens de votos válidos para os candidatos.
	 
\quebra	 

	% Média, mediana e moda 
	\item Faça um programa que calcule a média, a mediana e a moda dos valores de uma lista.

	\item Faça um programa que calcule a média, a mediana e a moda das chaves de um dicionário, no qual os valores associados às chaves indicam suas frequências.

	% decomposição em primos 	
	\item Faça um programa que exiba a decomposição em fatores primos de um inteiro. 
	
	Faça duas versões: exibindo os fatores repetidos e exibindo os fatores com expoentes, se forem repetidos. 
	
	% lista de divisores 
	\item Faça um programa que liste todos os divisores positivos de um inteiro.
	
	% números perfeitos
	\item Faça um programa que informe se um inteiro é um número perfeito: 
	é igual à metade da soma dos seus divisores.
	
	% números amigos 
	\item Faça um programa que informe se dois inteiros são números amigos: 
	a soma dos divisores de um é igual ao dobro do outro.
	
	Altere o programa para aceitar mais inteiros e verifique se formam um ciclo de números amigos:
	a soma dos divisores do primeiro é o dobro de outro, cuja soma dos divisores é o dobro de um terceiro.... 
	até que a soma dos divisores do último é o dobro do primeiro.

\quebra	
	
	\item Faça um programa que receba um inteiro e repetidamente aplique as operações
	\begin{equation*}
		\begin{cases}
			\frac{n}{2}	& \text{se $n$ for par}
			\\ 3n + 1	& \text{se $n$ for ímpar}
		\end{cases}
	\end{equation*}   
	até $n$ ser 1.
	
	Exiba todos os passos.
	
	Altere o programa para receber vários inteiros e, para cada um, mostrar seus passos.
	
	Crie um dicionário para não precisar recalcular as operações se repetir algum valor de $n$.
	
	\item Crie um programa que percorra um dicionário a partir de uma chave inicial. 
	
	Para cada chave, acesse seu valor e utilize-o como a próxima chave. 
	
	Se o valor de uma chave não for outra chave do dicionário, encerre o percorrimento.
	
	\item Crie um programa que percorra um dicionário a partir de uma chave inicial e identifique se há algum ciclo fechado.
	
	Utilize o algoritmo de detecção de ciclos de Floyd (lebre e tartaruga):
	percorra simultaneamente duas vezes, a cada passo da tartatuga, avance a lebre duas vezes.
	
	Quando a lebre atingir o final, sabe-se que não há ciclo.
	
	Se a lebre e a tartaruga estiverem na mesma posição em algum momento, um ciclo foi detectado.
	
	Altere o programa para identificar se há algum ciclo fechado, considerando início em todas as chaves.

\end{enumerate}